\subsection{Baryonic Uncertainties}
  \label{subsec:baryonic_uncertainties}

  The predictions of the effective framework depend explicitly on the observed baryonic mass distribution.
  Uncertainties in stellar mass-to-light ratios, gas content, and distance
  estimates therefore propagate into the inferred rotation curves.
  Such uncertainties are well documented
  and represent a generic limitation of baryon-based dynamical modeling~\cite{Lelli2016SPARC}.

  We have tested the sensitivity of the effective fits to reasonable variations in stellar mass normalization.
  While moderate shifts in mass-to-light ratio affect the detailed shape of the
  inner rotation curves, they do not eliminate the emergence of a low-acceleration
  saturation regime at large radii.
  In particular, the appearance of asymptotically flat rotation curves
  remains robust across the allowed baryonic parameter range, consistent with
  empirical scaling relations~\cite{McGaugh2016}.

  At cosmological scales, baryonic uncertainties primarily affect the mapping
  between large-scale structure and local observational samples.
  They do not remove the qualitative distinction between dense environments,
  which probe an unsaturated effective regime, and cosmic voids, which sample the
  saturated regime most strongly~\cite{Keenan2013KBC}.
